\chapter{Dehydration}

\section*{Definition}
A state of negative fluid balance resulting from excessive water loss (through skin, lungs, kidneys, or GI tract) or inadequate intake, leading to intracellular and extracellular volume depletion with potential electrolyte disturbances.

\section*{What Happens in Your Body}
Water loss exceeds intake, causing reduced total body water, increased plasma osmolality, and stimulation of osmoreceptors in the hypothalamus. Antidiuretic hormone (ADH) is released, increasing water reabsorption in renal collecting ducts. Thirst is stimulated. Continued loss without replacement leads to hypovolemia, reduced cardiac output, compensatory tachycardia, narrowing of blood vessels, and eventually hemodynamic instability and shock. Cellular dehydration impairs enzyme function and organ performance.

\section*{Causes}
\begin{itemize}
  \item Inadequate fluid intake (common in elderly, infants, dementia, during illness)
  \item Vomiting and diarrhea (gastroenteritis most common cause in children)
  \item Excessive sweating (exercise, fever, heat exposure)
  \item Polyuria (diabetes mellitus, diabetes insipidus, diuretic medications)
  \item Blood loss (hemorrhage)
  \item Burns (loss of skin barrier leading to fluid evaporation)
  \item Fever (increased insensible water loss)
  \item Alcohol intake (diuretic effect)
  \item Gastrointestinal losses (NG suction, fistulas, ostomy output)
  \item Reduced thirst sensation in elderly (presbyphagia, cognitive impairment)
\end{itemize}

\section*{When to See a Doctor}
See your doctor if:
\begin{itemize}
  \item Symptoms are severe or getting worse over time
  \item Severe dehydration signs: no urine output for 8 hours, lethargy or confusion, rapid heart rate, low blood pressure, sunken eyes, or inability to keep fluids down
  \item You have other concerning symptoms such as fever, weight loss, or bleeding
\end{itemize}

\section*{Related Symptoms}
\begin{itemize}
  \item Thirst
  \item Dark urine
  \item Fatigue
  \item Dizziness
  \item Dry mouth
\end{itemize}
\textbf{Important:} If this symptom persists, your doctor will check for serious underlying causes.

\section*{Self-Care / Home Management}
\begin{itemize}
  \item Rest and avoid activities that make it worse.
  \item Maintain adequate hydration and a balanced diet.
  \item Over-the-counter remedies may provide symptomatic relief (consult a pharmacist or doctor).
  \item Monitor symptoms and seek professional advice if they persist or worsen.
  \item Avoid self-medicating with prescription medications without medical guidance.
\end{itemize}

\section*{How Your Doctor Will Evaluate You}
\begin{itemize}
  \item A thorough history and physical examination are the first steps in evaluation.
  \item Laboratory tests (blood work, urinalysis) may be ordered based on clinical suspicion.
  \item Imaging studies (X-ray, ultrasound, CT, MRI) may be indicated when structural pathology is suspected.
  \item Specialist referral is arranged when the diagnosis remains uncertain or requires specific expertise.
\end{itemize}

\section*{Treatment Options}
\begin{itemize}
  \item Treatment targets the underlying cause identified through diagnostic evaluation.
  \item Supportive care and symptom management are provided while awaiting definitive therapy.
  \item Medications, lifestyle modifications, and physical therapies may be prescribed as appropriate.
  \item Follow-up ensures treatment effectiveness and allows timely adjustments to the care plan.
\end{itemize}

\clearpage